\section{Conclusion}
\label{sec:concl}

We have presented an extension of \GROOVE that enhances the tool's practical usefulness in incorporating data, based on user-defined record types and arbitrary operations. This work was inspired by a feature request.\footnote{Thanks to Reiko Heckel and Adam Machowczyk for their suggestions and feedback}

\subsection{Limitations}
\label{sec:limitations}

We list some practical limitations in the current \GROOVE implementation.
%
\begin{itemize}
\item \emph{User types and operations are only supported in the \javaA algebra.} In other words, the choice of algebras (\termA, \bigA, \javaA, \pointA) recalled in the introduction is currently not available for grammars with used-defined types. We expect that the extension to \termA and \pointA will be straightforward; however, to also be compatible with \bigA, the user types and operations would natively have to rely on \code{BigInteger} and \code{BigDecimal} rather than \intT and \doubleT.

\item \emph{User types are restricted to \recordT{}s.} This is a choice motivated by the fact that Java \recordT instances are inherently immutable and \recordT{}s have well-defined constructors. We expect that this will not be a limitation in practice.

\item \emph{User type fields are restricted to primitive sorts.} Though there is no inherent problem in allowing user types to have user-typed fields, cyclic dependencies would necessitate the introduction of \nullT values.

\item \emph{User type checking is weak.} The \GROOVE type system cannot statically distinguish different (Java) user types: they are all values of the single \userS sort. Run-time type errors result in the special $\bot_\userS$ error value (see \Cref{sec:system}).

\item \emph{Operator names must be distinct.} The current implementation does not allow overloading. In particular, this means that also field names must be distinct across \UserType classes.

\item \emph{Indeterminate operations cannot appear in quantified rules.} For instance, it is not possible to simultaneously assign random values to a set of nodes by universally quantifying the nodes; or for the example in this paper, neither of the rules in \Cref{fig:rules} can be universally quantified so as to apply to all \code{Person} nodes simultaneously.
%This is due to the fact that, in the current implementation, ``old'' rule matches of amalgamated rules cannot be reconstructed if they rely on past outcomes of an indeterminate operation, since those past outcomes are not always recorded.
Lifting this restriction will require a fairly major restructuring of \GROOVE's rule matching engine.
\end{itemize}


\subsection{Related work}
\label{sec:related}

%The integration of data attributes into graph transformation has been formalised and implemented in different ways over the years. With respect to formalisation, 
The seminal work of Ehrig et al.\ \cite{Ehrig-attributes} embeds the theory of attributed graphs (in which actually not just nodes but also edges can have attributes) and their transformation into algebra. \GROOVE's implementation is based on a variant of this embedding (see \cite{GROOVE-attributes}); for practical purposes, the main difference is that only nodes can be attributed and that attribute types are restricted to the four primitive sorts recalled in \Cref{sec:background}.

%Although basing graph attributes on algebra opens the door to a wealth of theory, it inherently limits data manipulation to deterministic operations, which is incompatible with randomness. We have only been able to resolve this by formalising our ``indeterminate'' operators as relations rather than functions.

Algebraic theory obviously allows much richer data types than just \GROOVE's primitive attributes. However, once users are allowed to program their own types and operations, the equational reasoning that is foundational to algebra is very hard to uphold. This is also seen in the tool \href{https://www.user.tu-berlin.de/o.runge/AGG/WWW/agg.html}{\AGG}\footnote{A more recent, open source snapshot of \AGG can be found \href{https://github.com/Rentenatus/oagg}{here}} (see \cite{AGG,AGG-2}): they, too, essentially import Java types into graphs, at the price of disabling some of \AGG's analysis capabilities, e.g., critical pair analysis.

Again based on the theory of algebraic node and edge attributes, Lambers and Orejas et al.\ \cite{Lambers-attributes} have developed a very powerful theory of symbolic reasoning, which in that paper is reported as having been implemented in a tool called \AutoGraph. (Unfortunately, the tool itself seems to be unavailable.)

\medskip\noindent A different way of integrating attributes and graphs has been both theoretically worked out and implemented by Plump et al.\ in the context of the Graph Programming tool \href{https://github.com/UoYCS-plasma/P-GP2}{GP} (see \cite{GP,GP-2}). Rather than aiming to be all-encompassing by allowing any algebraic type, they restrict data attributes to a compact, powerful but closed set of data types consisting of integers, strings and lists of those two, plus booleans.

\medskip\noindent On the other side of the spectrum, there are graph transformation tools that were conceived to be usable in practice, rather than to embrace theoretical concepts, and hence do not strive for a complete formalisation. Prime among these is \href{https://github.com/eclipse-henshin}{Henshin} (see, e.g., \cite{Henshin,Henshin-usability}), which actually shares many of \GROOVE's capabilities regarding state space exploration and model checking. Citing \cite{Henshin}, for data attributes they offer ``flexible attribute computations based on Java or JavaScript.''
% Another example is \href{https://grgen.de/}{GrGen.NET} (see \cite{GrGen,GrGen-NET}), in which data attributes either belong to one of five primitive types (the ones also supported by \GROOVE, except that \realS is split into \floatS and \doubleS), user-defined \enumS-types, plus the type \objectS that essentially corresponds to the top of the inheritance hierarchy and therefore can be instantiated by any user-defined type --- in other words, it is in principle quite similar to our \userS sort.
